\section*{Chapitre \ref{ch:m1:options-market-structure} --- Structure du marché des options}

\begin{solution}{exo:m1:options-market-structure:1}
C1 reçoit 10 et C2 reçoit 2~: les clients sont servis dans l'ordre chronologique, et l'ordre
est épuisé avant qu'aucun professionnel ne soit atteint.
\end{solution}

\begin{solution}{exo:m1:options-market-structure:2}
Clients~: 10 et 5, laissant 100. La part au prorata du \omterm{def:m1:what-a-trading-firm-does:market-maker}{teneur de marché} désigné
serait $100 \times 100/500 = 20$~; son droit est de 40. Les 60 restants sont
partagés entre les 400 lots de MM1, F et MM2~: $\lfloor 60 \times 100/400\rfloor
= 15$, $\lfloor 60 \times 50/400\rfloor = 7$, $\lfloor 60 \times 250/400\rfloor
= 37$~; le \omterm{def:m1:us-equity-market-structure:lots}{lot irrégulier} revient au plus ancien, MM1~: 16, 7, 37.
\end{solution}

\begin{solution}{exo:m1:options-market-structure:3}
Valides~: 2.43, 2.99 (centimes sous 3.00~\$), 3.05 et 12.10 (pas de cinq centimes à partir de
3.00~\$). Invalides~: 3.02 et 12.12. Autour de 3.10, la fourchette minimale vaut 0.05, soit
1.6~\% du prix.
\end{solution}

\begin{solution}{exo:m1:options-market-structure:4}
$311 \times 10^9 / 23\,400\,\text{s} = 13.3$ millions de messages par seconde en
moyenne. Pic~: 13.575 millions par 100 millisecondes, soit un débit de 135.8 millions par
seconde, 10.2 fois la moyenne. Bande passante~: 4.403 gigabits par 100 millisecondes
soit 44 gigabits par seconde pour un seul flux~; deux flux avec 10~\% pour
les retransmissions~: 97 gigabits par seconde.
\end{solution}

\begin{solution}{exo:m1:options-market-structure:5}
Personne~: l'initiateur exécute les 100 lots. Un répondant au stop~: 50. Trois
au stop~: 40, les répondants se partageant 60. Un à 2.49 pour 60 et deux à 2.50~:
les 60 lots améliorés vont à celui qui améliore le prix~; au stop, l'initiateur peut prendre
jusqu'à 40~\% de l'ordre initial, soit les 40 restants.
\end{solution}

\begin{solution}{exo:m1:options-market-structure:6}
Jambe par jambe~: payer 2.50, recevoir 0.90~: 1.60. Point médian~: $2.45 - 0.95 = 1.50$. Le
marché affiché du package est de 1.40 à 1.60, soit 0.20 de large. Le delta et
le vega du spread ne sont qu'une fraction de ceux de l'une ou l'autre jambe, si bien que
le risque pris par un \omterm{def:m1:what-a-trading-firm-does:market-maker}{teneur de marché} en le cotant est faible~; et les erreurs de ses deux évaluations de jambes se compensent largement. Un \omterm{def:m1:options-market-structure:complex}{carnet
d'ordres complexes} permet d'exprimer cela par un seul prix, à un ou deux centimes de 1.50.
\end{solution}

\begin{solution}{exo:m1:options-market-structure:7}
Avec un répondant en moyenne~: l'initiateur conserve 47.8~\% et 39.7~\% de l'ordre
est amélioré. Avec trois~: 12.7~\% et 77.5~\%. Le troisième répondant presque double la
part de l'ordre du client qui reçoit un meilleur prix, et prend trois quarts de ce que
conservait l'initiateur.
\end{solution}

\begin{solution}{exo:m1:options-market-structure:8}
Ici, la priorité temporelle n'existe pas entre professionnels. Tous les clients au
prix sont exécutés en premier~; puis les 100 lots sont partagés selon leur taille~: la part du cours acheteur de 5 lots
part est $\lfloor 100 \times 5/455\rfloor = 1$ lot, plus éventuellement le \omterm{def:m1:us-equity-market-structure:lots}{lot irrégulier}
issu de l'arrondi~: 2 lots au mieux, les \omterm{def:m1:what-a-trading-firm-does:market-maker}{teneurs de marché} 98. Pour être exécuté dans un tel
marché, un professionnel doit afficher de la taille, améliorer le prix ou répondre aux
enchères~; arriver en premier ne rapporte rien.
\end{solution}

\begin{solution}{pb:m1:options-market-structure:1}
\textbf{1.} 100, 50, 40, 40 (les 60 lots améliorés partent ailleurs et l'initiateur
prend les 40 restants, soit ses 40~\% de l'ordre initial) et 0.
\textbf{2.} $10 + 12.5 + 14 + 8 + 0 = 44.5$ lots.
\textbf{3.} Il vend à 2.50 ce qu'il évalue à 2.44~: 6 centimes par action, soit
6~\$ par contrat.
\textbf{4.} $44.5 \times 6 = \$267$.
\textbf{5.} Un centime sur 60 lots dans 20~\% des cas et sur 100 lots dans 10~\%~:
$0.2 \times 60 + 0.1 \times 100 = 22$ lots améliorés de 1~\$ chacun, soit 0.22~\$ par
contrat en moyenne sur l'ordre.
\textbf{6.} 2.50, le cours vendeur affiché, si 100 lots sont là. Le \omterm{def:m1:retail-flow-and-wholesaling:wholesaler}{wholesaler}
ne gagnerait rien sur la transaction elle-même.
\textbf{7.} $267 - 100 \times 0.30 = \$237$.
\textbf{8.} Le client paie au plus 2.48~: deux centimes d'amélioration pour
certain. La marge du \omterm{def:m1:retail-flow-and-wholesaling:wholesaler}{wholesaler} tombe à 4 centimes par action. Les répondants doivent désormais
proposer 2.47 pour améliorer, ce que peu feront à une valeur de 2.44 à 2.45~: le
\omterm{def:m1:retail-flow-and-wholesaling:wholesaler}{wholesaler} conserve une part plus grande d'une marge plus mince. Le prix stop est le
choix du \omterm{def:m1:retail-flow-and-wholesaling:wholesaler}{wholesaler} et son principal instrument.
\textbf{9.} À 2.50 et 2.49 (marges de 5 et 4 centimes)~; à 2.46, il n'a qu'un centime.
Être à temps signifie recevoir le message d'enchère, reconnaître la série,
recalculer le prix à partir du cours actuel de l'action et de la volatilité, vérifier le risque et
répondre, le tout bien en deçà de 100 millisecondes moins deux allers-retours réseau~: un
moteur de prix automatisé et colocalisé.
\textbf{10.} Pas comme une cotation au repos en concurrence avec les réponses à conditions égales~:
les clients prioritaires du carnet sont honorés, les cotations professionnelles ne sont pas
garanties d'une part d'un ordre apporté à l'enchère déjà
apparié. Pour y participer, il doit envoyer une réponse.
\textbf{11.} Répondre révèle son prix pour cette série à l'initiateur
et à tout participant qui voit l'enchère, cent fois par seconde~; et
les ordres apportés aux enchères sont sélectionnés. Il peut néanmoins répondre, car
les flux soumis aux enchères sont les meilleurs qui soient.
\textbf{12.} Les ordres de professionnels, les grands ordres institutionnels exécutés par
des algorithmes, et les ordres de firmes réagissant aux nouvelles~: le flux qui atteint
les cotations affichées non accompagné est, en moyenne, mieux informé.
\textbf{13.} Plus large. Une cotation affichée est une option offerte aux participants les plus rapides et
les mieux informés~; elle doit être valorisée pour eux.
\textbf{14.} C'est la segmentation décrite au \cref{ch:m1:retail-flow-and-wholesaling}
mise en œuvre sur une bourse~: les flux non informés sont isolés, cotés au plus serré et partagés
entre ceux qui y ont un droit~; la cotation affichée est laissée pour servir
tous les autres.
\textbf{15.} Pour~: plusieurs firmes bien capitalisées voient toutes les enchères et
répondent par machine~; un dixième de seconde est long pour elles~; l'amélioration
a bien lieu. Contre~: seules les firmes colocalisées peuvent répondre~; l'initiateur choisit le
stop et conserve jusqu'à la moitié quoi qu'il arrive~; un répondant qui se contente d'égaler
obtient une part minoritaire, si bien que l'incitation à améliorer est faible à moins qu'il ne compte
remporter l'ordre entier.
\textbf{16.} Celui qui négocie contre les cotations des \omterm{def:m1:what-a-trading-firm-does:market-maker}{teneurs de marché}. Ces frais sont un
coût de la tenue de marché sur les flux clients, récupéré dans la fourchette~; les clients
en reçoivent une partie en retour, grâce à la tarification sans commission de leurs courtiers.
\textbf{17.} Pour servir les deux clientèles à la fois, \omterm{def:m1:what-a-trading-firm-does:market-maker}{teneurs de marché} sur une bourse
et preneurs de liquidité sur une autre, pour obtenir davantage de droits et de revenus
de données de marché et de connectivité, et parce que chaque licence de bourse
constitue un emplacement distinct dans la table de routage de chaque courtier.
\textbf{18.} Moins à améliorer~: avec une fourchette d'un ou deux centimes, le stop est
déjà proche de la fair value, les enchères amélioreraient moins, et les marges des \omterm{def:m1:retail-flow-and-wholesaling:wholesaler}{wholesalers}
diminueraient, de même que la taille affichée.
\textbf{19.} \textbf{44.5 lots sur les 100.}
\textbf{20.} Quelques dollars par contrat pour la firme qui l'exécute, ce qui explique pourquoi
les courtiers sont rémunérés pour cela, les bourses construisent des mécanismes autour, et le
marché affiché est ce qui reste.
\end{solution}

\begin{solution}{iq:m1:options-market-structure:1}
Une chambre de compensation et de nombreuses bourses, comme en actions, mais~: des centaines de
fois plus d'instruments, presque toute la liquidité provenant des cotations des \omterm{def:m1:what-a-trading-firm-does:market-maker}{teneurs de marché}
plutôt que des ordres naturels~; \omterm{def:m1:options-market-structure:customer}{priorité client} et \omterm{def:m1:matching-algorithms-and-implied-spreads:prorata}{allocation au prorata}, avec des droits des
\omterm{def:m1:what-a-trading-firm-does:market-maker}{teneurs de marché} sur de nombreuses bourses~; des frais qui dépendent de la capacité du
trader~; des \omterm{def:m1:options-market-structure:auction}{enchères d'amélioration de prix} comme principal canal pour le flux de détail, au lieu de
l'\omterm{def:m1:retail-flow-and-wholesaling:wholesaler}{internalisation} hors bourse~; des \omterm{def:m1:options-market-structure:complex}{carnets d'ordres complexes}~; et un flux consolidé de
centaines de milliards de messages par jour.
\iqlookfor{cotation pilotée, règles dépendant de la capacité, enchères.}
\end{solution}

\begin{solution}{iq:m1:options-market-structure:2}
Les ordres au repos des clients particuliers au meilleur prix sont exécutés avant ceux de tout
professionnel, quel que soit l'instant d'arrivée. Les bourses l'accordent pour attirer les ordres
des clients, que les professionnels apprécient parce qu'ils sont non informés~; le flux
client attire ensuite les \omterm{def:m1:what-a-trading-firm-does:market-maker}{teneurs de marché}, dont les cotations font le
marché de la bourse.
\iqlookfor{la logique de plateforme à deux versants.}
\end{solution}

\begin{solution}{iq:m1:options-market-structure:3}
Initiateur~: choisir les ordres clients à mettre aux enchères et le stop~; plus le stop est bas
à l'achat, meilleur est le prix pour le client et plus mince la marge,
mais plus grande la part conservée. Soumettre la paire~; la bourse garantit au client le stop. Résultat~:
les réponses au meilleur prix sont retenues en premier~; au stop, les clients présents au carnet,
puis jusqu'à 50~\% ou 40~\% pour l'initiateur, puis le reste. Répondeur~: recevoir le message,
valoriser la série immédiatement, décider entre égaler (une part au prorata de ce que
l'initiateur laisse) et améliorer (droit de premier rang), et répondre
dans le délai imparti.
\iqlookfor{le prix du stop comme variable stratégique~; égaler ou améliorer.}
\end{solution}

\begin{solution}{iq:m1:options-market-structure:4}
Volume et rafales~: des centaines de milliards de messages par jour, avec des pics dix fois
supérieurs à la moyenne en 100 millisecondes, des dizaines de gigabits par seconde par flux, deux
flux à arbitrer, détection des trous et retransmission. Filtrer aussi tôt que possible (la plupart
des séries ne concernent aucune stratégie donnée), idéalement dans le matériel ou avec des
récepteurs à contournement du noyau~; correspondance des symboles pour plus d'un
million de séries qui changent chaque jour~; construction des carnets par série et par bourse~;
et politiques de condensation pour les consommateurs lents qui ne réordonnent ni n'inventent
jamais de prix.
\iqlookfor{conception pour les pics par rapport à la moyenne, filtrage précoce, arbitrage de ligne.}
\end{solution}

\begin{solution}{iq:m1:options-market-structure:5}
Toutes les cotations d'un \omterm{def:m1:what-a-trading-firm-does:market-maker}{teneur de marché} dans une classe reflètent une seule opinion sur le
sous-jacent~; si cette opinion est périmée, un trader informé peut toucher des centaines
d'entre elles avant que le teneur ne reprice. La protection limite les dégâts en annulant tout
ce qui reste après un nombre d'exécutions dans une fenêtre donnée. Le coût~: la taille affichée
est en partie illusoire, puisqu'elle disparaît après les premières exécutions, de sorte
qu'un preneur de taille voit les cotations s'évanouir~; et les mécanismes eux-mêmes peuvent retirer
brutalement la liquidité sur des marchés rapides.
\iqlookfor{corrélation des cotations comme raison~; évanouissement comme coût.}
\end{solution}

\begin{solution}{iq:m1:options-market-structure:6}
Le flux le moins informé est détourné vers les enchères et alloué selon les droits, de sorte que
ce qui atteint ma cotation affichée sans accompagnement est sélectionné adversement~: professionnels,
algorithmes, nouvelles. Je devrais coter le carnet plus large que ma fourchette
pour le flux moyen, le dimensionner pour le flux toxique, le protéger par un repricing rapide
piloté par le sous-jacent et des limites de risque, et me battre pour le flux bénin là où il
circule réellement~: répondre aux enchères, obtenir un rôle de \omterm{def:m1:what-a-trading-firm-does:market-maker}{teneur de marché} désigné lorsque
le droit rémunère les obligations, ou négocier le flux dirigé.
Mesurer séparément les markouts de chaque canal.
\iqlookfor{\omterm{def:m1:what-a-trading-firm-does:adverse-selection}{sélection adverse} du flux résiduel coté~; aller là où se
trouve le flux.}
\end{solution}
